\subsection{Assignment 15 -- Distribuzione mensile degli stream per regione e il totale della nazionale}

L’obiettivo dell'analisi è la rappresentazione, su base mensile, di una struttura ordinata che
elenca tutte le regioni appartenenti a ciascun paese, affiancata da una riga
di sintesi per il totale nazionale e del relativo numero di
stream osservati. Per rendere esplicita e leggibile tale struttura nel risultato
finale, sono stati definiti due membri calcolati,
\textit{Paese\_Col} e \textit{Regione\_Col}, che consentono di visualizzare
 le etichette testuali del paese e della regione nel set di risultato.

Questo approccio permette di gestire correttamente la riga del totale
nazionale sostituendo la regione con una voce aggregata (\textit{Total}), senza
la necessità di creare nuovi elementi all’interno della gerarchia geografica del
cubo.

La costruzione dell’insieme completo delle righe è realizzata mediante la funzione
\textit{GENERATE}, che per ciascun paese produce l’elenco delle regioni
associate, ottenute tramite la funzione \textit{DESCENDANTS} limitata al livello
\textit{Region}, e successivamente una riga aggiuntiva corrispondente al totale
nazionale. La combinazione di tale struttura con il livello \textit{Month} della
gerarchia Date consente di replicare, per ogni mese, la sequenza
\emph{Paese–Regioni–Totale}.

La clausola \textit{FILTER} è utilizzata per rimuovere i membri vuoti o
privi di nome, mentre l’uso di \textit{NON EMPTY} consente di escludere
automaticamente le combinazioni per cui il numero di stream risulta assente.
Infine, l’ordinamento del risultato è ottenuto concatenando mese, paese e regione,
in modo da garantire una presentazione interpretabile dei dati.
